% Generated draft: requires mathematical and source review.
\documentclass[11pt]{article}
\usepackage[margin=25mm]{geometry}
\usepackage{fontspec}
\setmainfont{Latin Modern Roman}
\usepackage{amsmath,amssymb,mathtools,unicode-math}
\setmathfont{Latin Modern Math}
\usepackage{enumitem,needspace,xurl,hyperref}
\setlength{\parindent}{0pt}
\setlength{\parskip}{5pt}
\setlength{\emergencystretch}{3em}
\newcommand{\N}{\mathbb{N}}
\newcommand{\Z}{\mathbb{Z}}
\newcommand{\Q}{\mathbb{Q}}
\newcommand{\R}{\mathbb{R}}
\newcommand{\C}{\mathbb{C}}
\newcommand{\F}{\mathbb{F}}
\newcommand{\A}{\mathbb{A}}
\newcommand{\PP}{\mathbb{P}}
\newcommand{\E}{\mathbb{E}}
\newcommand{\eps}{\varepsilon}
\newcommand{\dd}{\,\mathrm{d}}
\begin{document}
\section*{UnsolvedMath ID: 30005921}
\textbf{Source problem code:} catalogue entry
\subsection{Definitions and mathematical statement}
The supplied record is introductory setup, not a recoverable standalone conjecture. It refers to conjectured generalizations of equations (3) and (4), whose statements are absent. Its supplied verification note identifies the intended later targets as Conjectures 10 and 12, preserved in subsequent records; neither conjecture is included here.

The recoverable mathematical content is an attributed lemma concerning a polynomial \(p(x_1,\ldots,x_d)\) in noncommuting indeterminates, normalized by \(p(0)=1\). For a matrix tuple \(X=(X_1,\ldots,X_d)\), evaluation means substituting the matrices into each word of \(p\), with the constant term evaluated as the identity matrix. The extraction does not explicitly specify the coefficient field, the matrix field, or the definition of \(\|\cdot\|_{\mathrm{row}}\). A usual complex-matrix convention is
\[
\|X\|_{\mathrm{row}}=\left\|\sum_{j=1}^{d}X_jX_j^*\right\|_{\mathrm{op}}^{1/2},
\]
where \(^*\) denotes conjugate transpose and \(\|\cdot\|_{\mathrm{op}}\) is the operator norm for the Euclidean norm. This convention is recorded for interpretation only; its applicability to the source must be checked.

With those unresolved conventions understood, the passage labeled \emph{Lemma 9} asserts the equivalence of the following conditions:
\begin{enumerate}
\item For every positive integer \(n\) and every tuple \(X=(X_1,\ldots,X_d)\) of \(n\times n\) matrices satisfying \(\|X\|_{\mathrm{row}}\leq 1\), one has
\[
\det p(X_1,\ldots,X_d)\ne 0.
\]
\item There exist a positive integer \(m\) and a tuple \(A=(A_1,\ldots,A_d)\) of \(m\times m\) matrices satisfying \(\|A\|_{\mathrm{row}}<1\) such that, for every positive integer \(n\) and every tuple \(X=(X_1,\ldots,X_d)\) of \(n\times n\) matrices, without a norm restriction,
\[
\det p(X_1,\ldots,X_d)
=
\det\!\left(I_m\otimes I_n-\sum_{j=1}^{d}A_j\otimes X_j\right).
\]
Here \(I_k\) denotes the \(k\times k\) identity matrix and \(\otimes\) is the Kronecker product. The summation in this determinant is an explicit reconstruction of a damaged summation marker in the extraction.
\end{enumerate}
This is a qualified transcription of an attributed lemma, not an independently verified theorem or a replacement for the missing conjectures.

The extraction also interleaves damaged equations (7) and (8). Equation (7) concerns coefficients \(c_\alpha(N)\), with multi-index \(\alpha=(\alpha_1,\ldots,\alpha_d)\), expressed using binomial and multinomial coefficients; the placement of factors and exponents cannot be recovered reliably. Equation (8) describes an \(N\to\infty\) integral over \(U(N)^d\) involving determinants built from \(\sum_j z_jU_j\) and \(\sum_j w_jU_j^*\), with a displayed right-hand side \((1-\langle z,w\rangle)^{-1}\). Its hypotheses, complex conjugations, and integration normalization are insufficiently clear to state a precise identity. These fragments supply no recoverable additional conjecture.
\subsection{Short English statement}
This record introduces a proposed notion of noncommutative stability through an attributed equivalence between determinant nonvanishing on a closed row-norm ball and a determinantal representation with strictly contractive coefficient matrices. The actual conjectures are missing, and several conventions and interleaved formulas require source review. The supplied open label concerns an incomplete record; it does not establish that Lemma 9 is an open problem.
\subsection{Sources}
\begin{itemize}
\item Local UnsolvedMath registry, permanent ID \(30005921\), source code \texttt{catalogue entry}, title \emph{Setup for Noncommutative Stability Conjectures}; supplied complete source record. No individual proposer is identified.
\item \emph{Non-commutative Function Theory and Free Probability} (2024), Oberwolfach Reports, pp.\ 1227--1284, DOI \texttt{10.4171/owr/2024/22}. Supplied report reference and literature-source label: OWR 2024/22. \url{https://doi.org/10.4171/owr/2024/22}
\item Supplied literature-source label: Volcic. \url{https://arxiv.org/abs/1807.05645}. The record supplies no further bibliographic details or precise theorem citation for this reference.
\item Literature assessment dated exactly \texttt{2026-08-22}: supplied dated curation, not independently verified. It labels the record conservatively open because the extraction is an incomplete transition to a conjectured stability notion, and states that known determinantal-representation theorems do not identify or settle the later conjectures. The background expressly calls for independent verification of time-sensitive status and cited claims.
\item The supplied statement status is \texttt{unrecoverable}. Its verification note reports that the official report confirms introductory setup for Conjectures 10 and 12 in subsequent records. This is a supplied verification claim; no independent source consultation is asserted here.
\end{itemize}
\end{document}
